\begin{textsample}{NEG Dim 1 – human – Score -7.00 – rj/rj\_ef\_linguagens\_ed\_fis\_2\_p1.txt}  \label{ex:f1_neg_005}
A Educação Física no Ensino Fundamental é o componente curricular que \textbf{tem} como objeto de estudo a cultura \textbf{corporal} de \textbf{movimento} e \textbf{tem} como função \textbf{possibilitar} aos estudantes o acesso a práticas diversificadas na Área de Linguagens, com \textbf{significados} e regionalidades que imprimem características únicas ao \textbf{movimento} e às expressões da cultura \textbf{corporal}, além de trazerem uma característica singular de expressão de \textbf{determinado} gesto, pois sendo culturalmente definido, carrega consigo, e por isso demonstra, uma \textbf{maneira} única de ser expresso, de se apresentar ao mundo, de se comunicar, destacando, assim, a importância da Educação Física como um componente da Área de Linguagens e suas Tecnologias ( DAÓLIO, 2010 ; NEIRA, 2009 ; NEIRA \& SOUZA JÚNIOR, 2016 ).

% matched lemmas: corporal, determinar, maneira, movimento, possibilitar, significado, ter
\end{textsample}
